
Immigration remains one of the most hotly discussed topics in economics and the political sphere. Actors from across the political and ideological spectrum share conflicting opinions on the effects of migration. As such, they advocate for widely different levels of immigration in North America. Ranging from completely open borders to mass deportations, and everything in between. An important piece of the puzzle for creating optimal immigration policy is understanding how immigrants differ from natives. I study this question in the dimension of managing businesses. I am interested in looking and at immigrant versus native owned businesses and understanding how the two groups differ in characteristics. I conduct this research in the setting of the United States, which provides a uniquely beneficial environment for a multitude of reasons. It is a unique country with large and sustained waves of immigration throughout its history. Even in 2025, about 15 \% of the US population were immigrants. It's also a large country both in economic activity and geographical area, which provides unique historical variation one can leverage for identification purposes. 

In this paper, I attempt to look at immigration through a model of firm size distribution and infer the underlying managerial talent of immigrants. Secondly, I utilize an instrumental variable approach to extract causal estimates for the effects of immigration. The rest of the paper is organized as follows: Section \ref{sec:lit} situates my research among existing work in the field, and Section \ref{sec:data} discusses the data sources. Section \ref{sec:model} explains my theoretical model, and Section \ref{sec:cross_section} provides the model prediction and evidence to back them. Section \ref{sec:results} provides the IV results, Section \ref{sec:discussion} reconciles the theoretical and empirical work, and Section \ref{sec:conclusion} concludes.